\prefacesection{Abstract}
\begin{center}
	\setstretch{2}
	\MakeUppercase\ThesisTitle
\end{center}
\doublespacing
\begin{center}
	\setstretch{2}
	\AuthorName
\end{center}
\setstretch{2}
% do not change anything above
\setlength{\parindent}{0in}
% the first paragraph of your abstract begins follows with 0 inch indent
% set the indent for the rest paragraphs to 0.5 inch
% \setlength{\parindent}{0.5in}

Inconsistent development environments pose substantial friction for software engineers and development teams, frequently consuming days of onboarding time and shifting engineering focus from feature delivery to environment debugging. In computer science education, this challenge is acutely amplified. When students must imperatively configure their personal machines for diverse coursework, valuable lecture and lab hours are often sacrificed to troubleshooting dependency conflicts. The diversity of student hardware, operating systems, and local configuration states turns instructors and teaching assistants into ad-hoc technical support providers, creating severe barriers to student engagement and learning.

% the first paragraph of your abstract ends here
% set the indent for the rest paragraphs to 0.5 inch
\setlength{\parindent}{0.5in}
% the rest paragraphs should begin as follows

This thesis proposes an automated framework that leverages agentic large language models (LLMs) to synthesize reproducible development environments using Nix Flakes from instructors' natural language requirements and unconfigured source code repositories. Because environment synthesis is a recurrent task across semesters and course suites, we designed and implemented a specialized custom agent skill that equips the agent with domain-specific flake templates, structured metadata, and automated validation heuristics to ensure high behavioral consistency. To prevent package hallucination within the massive Nixpkgs ecosystem, the specialized agent queries a Model Context Protocol (MCP) server to retrieve real-time package metadata and configuration schemas directly from NixOS search indexes. Synthesized flakes are verified through a closed-loop deterministic evaluation pipeline using native Nix tooling, enabling the agent to autonomously interpret compiler and evaluation error traces to self-correct invalid specifications. Furthermore, this study presents an empirical benchmarking of computational overhead, demonstrating that declarative Nix Flake environments offer substantial memory savings, disk deduplication, and reduced startup latency compared to conventional OCI DevContainers on resource-constrained student laptops.
